\documentclass{report}
\usepackage{amsmath,amssymb}
\setlength{\parindent}{0mm}
\setlength{\parskip}{1em}
\begin{document}
\begin{center}
\rule{6in}{1pt} \
{\large Verena Kuhlemann \\
{\bf Improving the communication pattern in mat-vec operations for large scale-free graphs by disaggregation}}

Mathematics & Computer Science \\ Emory University \\ Suite W401 \\ 400 Dowman Drive \\ Atlanta \\ GA 30322
\\
{\tt vkuhlem@emory.edu}\\
Panayot S. Vassilevski\end{center}

Matrix vector multiplication (mat-vec) is the key operation in any
Krylov-subspace iteration method. We are interested in Krylov methods
applied to problems associated
with graph Laplacians arising from large scale-free graphs. Computations
with graphs of this type on parallel distributed-memory computers are
challenging. This is due to the fact that scale-free graphs have a degree
distribution that follows a power-law and currently available graph
partitioners are not efficient for such an irregular degree distribution.
The lack of a good partitioning leads to excessive inter-processor
communication requirements during every mat-vec. We present an approach
to alleviate this problem based on embedding the original irregular graph
into a more regular one by disaggregating (splitting up) vertices in the
original graph. The mat-vec operations for the original graph are
performed via a triple matrix product involving the embedding graph. Even
though the latter graph is larger, we are able to decrease the
communication requirements considerably and improve the performance of
mat-vec.

This work was performed under the auspices of the U.S. Department of
Energy by Lawrence Livermore National Laboratory under Contract DE-AC52-07NA27344.


\end{document}
